\documentclass[11pt,letterpaper,reqno]{amsart}
\usepackage[T1]{fontenc}
\usepackage{lmodern,microtype}
\usepackage[letterpaper,margin=1in,headheight=15pt]{geometry}
\usepackage{amsmath,amssymb,enumitem,xcolor}
\definecolor{manuscriptred}{rgb}{0.6,0.1,0.1}
\definecolor{manuscriptblue}{rgb}{0.1,0.3,0.5}
\definecolor{manuscriptmagenta}{HTML}{60445C}
\usepackage[colorlinks=true,linkcolor=manuscriptred,citecolor=manuscriptblue,urlcolor=manuscriptmagenta]{hyperref}
\hypersetup{pdftitle={QED3 revisions: 2026-10-23},pdfauthor={}}
\setcounter{tocdepth}{1}
\title[QED3 revisions: 2026-10-23]{QED3 revisions: 2026-10-23}
\date{}
\begin{document}
\begin{abstract}
Minor corrections and local clarifications to the comprehensive QED3
manuscript and the corresponding companion passages, dated 23 October 2026.
The main conclusions are unchanged. The local fold-proof correction
restores the supplied argument already present in the reflection-positivity
companion; no new construction or estimate is introduced.
\end{abstract}
\maketitle
\tableofcontents

\section{Corrections}
The following list covers audit items QED-0001--QED-0042. Corrections
already present in a companion were retained. Matching TeX sources and
PDFs were refreshed in the repository and Dropbox collaboration copy.
\begin{enumerate}[label=\arabic*.,leftmargin=*,itemsep=4pt]
\item Clarified the abstract's weak-charge and ergodic-component scope,
the lattice-unit side length, charge-uniformity of the cumulant remainder,
and the one-parameter translation group used in the clustering proof
(1--4).
\item Allowed the notation table to break between rows, kept the
two-line continuity estimate together, and replaced inaccurate
``first-page statement'' references by theorem references (5, 34, 42).
\item Specified Euclidean face area, described the hopping map as a
phase-weighted site permutation, and stated the bound for each Grassmann
coefficient (6, 7, 14).
\item Made spelling consistent within each manuscript; corrected
punctuation and subject--verb agreement (8, 32, 37).
\item Assigned upper-slab coefficient support to $S_+$ and explicitly
retained the generator support and coefficient-variable dependence in
$\mathcal G^+$ (9, 10).
\item Identified the integer Fourier coefficients in the cut-link
integral, compared Berezin integration functionals consistently, and
defined the normalized Gaussian at first use (11, 12, 16).
\item Based each folded cyclic trace at its own slab boundary, including
half-integer-width applications. The trace identity and chessboard
estimate are unchanged (13).
\item Fixed the dual-tree root, removed an undefined cutoff exponent,
and specified the unit-lattice block and its exact references (15, 17, 26).
\item Distinguished analytic bilinear pairings from auxiliary Hermitian
Fourier products, made the conjugated connection explicit, and separated
operator and trace norms from weighted block norms (18, 29, 36).
\item Distinguished thick layers from chambers, photon overlap
components from fermion face components, and localization-family indices
from torus indices. Stated the common incoming data required in primitive
comparisons (19, 30, 31, 33).
\item Retained the scalar coefficient of a constant polynomial, restored
the factor $2\pi/q$ in the local lift, and specified unit mirror normals
(20--22).
\item Corrected the Bochner-integrability explanation, displayed the
spectral projections in the finite-stage identity, and used absolute
coordinate differences for the positive displacement (23--25).
\item Named the chosen real summand rather than the full real part.
Renamed the scalar polymer-sum constant $K_{\mathrm{pol}}$, supplied its
definition and reference, and clarified the individual counterterm
bounds before summation; operator and compact-set uses of $K_0$ remain
unchanged (27, 28).
\item Restricted cutoff independence to the stop offset and physical
stop quantities, specified nonnegative Gaussian-tail thresholds,
distinguished basis momenta from twist-shifted symbol arguments, restored
the mean-zero transverse field domain, and cited the selected coefficient
subsequence (35, 38--41).
\end{enumerate}
The refreshed companion exports also retain the project's current muted
link colors. Original source inputs and historical copies are unchanged.
\end{document}
